\documentclass[11pt,a4paper]{article}
\usepackage[margin=25mm]{geometry}
\usepackage{amsmath,amssymb,amsthm,hyperref}
\setlength{\emergencystretch}{2em}
\newtheorem{theorem}{Theorem}
\newcommand{\gen}[1]{\langle#1\rangle}
\title{Highest rank for the natural alternating groups}
\author{Human Compression project: principal recovery of L044}
\date{}
\begin{document}\maketitle
The result recovered here is Theorem~1.1 of
Cameron--Fernandes--Leemans--Mixer, \emph{Highest rank of a polytope
for $A_n$}, arXiv:1605.09173v1 (2016). All component proofs and
their necessary finite certificates are retained. This paper checks
their complete induction and source interfaces; its length alone
does not measure the written replacement.

A string C-marking of rank $r$ is an ordered family of nonidentity
involutions $\rho_0,\ldots,\rho_{r-1}$ generating the group, with
$\rho_i\rho_j=\rho_j\rho_i$ whenever $|i-j|>1$, and satisfying
\[
 \gen{\rho_i:i\in I}\cap\gen{\rho_i:i\in J}
       =\gen{\rho_i:i\in I\cap J}
       \quad(I,J\subseteq\{0,\ldots,r-1\}).
\]
The standard correspondence between such markings and abstract
regular polytopes is acquired. We use the source's rank-zero
convention when no generating string C-marking exists.

\begin{theorem}\label{principal}
For $n\ge3$, the highest rank of a string C-marking generating $A_n$ is
zero for $n=3,4,6,7,8$, three for $n=5$, four for $n=9$,
five for $n=10$, six for $n=11$, and
\[
                         \lfloor(n-1)/2\rfloor
                         \qquad(n\ge12).
\]
\end{theorem}
\begin{proof}
The complete acquired finite maxima settle every degree from
three through fourteen; the values at twelve, thirteen and
fourteen are $5,6,6$. The complete L038 existence theorem
supplies a generating string C-marking of rank
$\lfloor(n-1)/2\rfloor$ for every $n\ge12$. Its retained
earlier constructions supply existence and the intersection
property only: their maximality is not invoked. Thus this lower
bound does not use the upper bound being recovered here.

It remains to prove the upper bound for $n\ge15$ by strong
induction. Assume it in every degree $12\le q<n$, retaining
the exact smaller finite maxima. Let
$A_n=\gen{\rho_0,\ldots,\rho_{r-1}}$ be any string
C-marking. Its noncommutation diagram is connected. Otherwise
the subgroups on its components commute, intersect trivially
by IP and give a nontrivial direct decomposition of the simple
group $A_n$. Each maximal omission
$G_i=\gen{\rho_j:j\ne i}$ is proper by independence.

If some $G_i$ is primitive, the complete primitive reduction
bounds its rank by $\lfloor(n-3)/2\rfloor$. Hence
\[
             r\le1+\lfloor(n-3)/2\rfloor
                   =\lfloor(n-1)/2\rfloor.
\]
Suppose next that a maximal omission is transitive imprimitive.
Choose a maximal nontrivial block system, with $m$ blocks of
size $k$, so $n=km$. If $k,m>2$, the complete large-block
theorem gives the required bound. If $k=2$, the complete
pair-block theorem gives it, with its smaller-degree premises
supplied by the current induction. Therefore, under a putative
failure of the bound, every maximal block system of every
transitive omission has exactly two blocks.

If a double omission $G_{i,j}$ inside one of these transitive
omissions is transitive, the complete binary-coordinate theorem
and its exact degree-sixteen certificate exclude the configuration.
Its hypothesis that every transitive maximal omission preserves
a binary partition has just been established; it is not assumed
for an arbitrary imprimitive action. Since binary partitions
have equal blocks, $n$ is even, and the large-degree application
has $n\ge16$, as required. If instead
all double omissions inside a transitive $G_i$ are intransitive,
the complete binary-omission assembly gives the desired bound.
Its induction hypothesis concerns only degrees $12\le q<n$.
Thus every transitive maximal omission is excluded under a
putative rank failure.

It remains that every $G_i$ is intransitive. An actual edge of
$\rho_i$ is eligible when it joins different $G_i$-orbits.
Every label has such an edge: otherwise adjoining $\rho_i$
would not make its omission transitive. If some label has exactly
one eligible edge, the complete singleton and unique-crossing
assembly gives $r\le\lfloor(n-1)/2\rfloor$. Its strict
smaller-degree hypotheses are supplied by the induction assumption,
not by Theorem~\ref{principal} at degree $n$.

Otherwise each label has at least two eligible edges. Choose two
per label. The complete paired count gives $2r\le n$. In the
equality case, the acquired normal form and connectedness theorem,
followed by the complete paired-ladder proof, produce a partition
into pairs preserved by every actual generator. Since $n\ge15$
and equality makes $n$ even, the ladder theorem's condition
$n=2r>8$ holds. This contradicts the natural primitivity of
$A_n$. Therefore $2r<n$, which is exactly the required bound.
Every possibility for the maximal omissions has been exhausted,
and the induction is complete.
\end{proof}

\paragraph{Attribution and retained inputs.}
Theorem~\ref{principal} is the historical result of the four source
authors. Their paper attributes the small maxima to
Fernandes--Leemans--Mixer (2012), and records the large-degree
bound as the solution of their earlier conjecture. No new solution
of that already settled conjecture is claimed here. The lower
existence constructions remain attributed through the complete
L038 recovery. Classical inputs such as Whiston's generator bound,
Mar\'oti's primitive-group theorem and Conder's order bound are
retained where their actual component proofs use them. Required
finite computations and the completeness of the relevant GAP
libraries remain declared premises, with exact execution records.

The present recovery includes new intermediate orbit, support,
parity, pairing and centralizer arguments, identified separately
in the component manuscripts; no priority is claimed for their
elementary ingredients. The literal arbitrary-independent claim
in source Proposition~3.3 is false. It is replaced only by the
complete corrected theorem for independent involutory strings,
with its explicit counterexample and bounded certificates retained.
The finer individual ranks in the source's Table~1 are not premises
or additional recovered claims of this principal-theorem contract.
The mathematical assembly is distinct from publication approval,
formal verification and a numerical compression audit.
\end{document}
